\subsection{CONSTRUCTION OF INVARIANT DIFFERENTIAL FORMS}

Lemma 2. Let \(\mathbf{G}\) be a Lie group, \(\mathrm{U}\) a symmetric open neighbourhood of \(e\) in \(\mathbf{G}\) , \(\mathrm{E}\) a complete normable space and \(\phi: \mathrm{U}^2 \to \mathrm{E}\) an analytic mapping. For all \(g \in \mathrm{U}\) , let \(\omega_g\) be the differential at the point \(g\) of the mapping \(h \mapsto \phi(g^{-1}h)\) . Then \(\omega\) is the restriction to \(\mathrm{U}\) of the left invariant differential form on \(\mathbf{G}\) whose value at \(e\) is \(d_e\phi\) .

Clearly \(\omega_{e} = d_{e}\phi\) . For all \(g \in \mathrm{U}\) and all \(t \in \mathrm{T}_{e}(\mathrm{G})\) ,

\[
\begin{array}{r l} \langle \omega_ {g}, \mathrm{T} _ {e} (\gamma (g)) t \rangle & = \langle d _ {g} (\phi \circ \gamma (g) ^ {- 1}), \mathrm{T} _ {e} (\gamma (g)) t \rangle \\ & = \langle d _ {e} \phi \circ \mathrm{T} _ {g} (\gamma (g) ^ {- 1}), \mathrm{T} _ {e} (\gamma (g)) t \rangle = \langle d _ {e} \phi , t \rangle \end{array}
\]

and hence \(\omega_{g}\) is derived from \(d_{e}\phi\) by \(\mathrm{T}_{e}(\gamma(g))\) .

PROPOSITION 52. Let \(n\) be an integer \(>0\) , \(G\) an \(n\) -dimensional Lie group, \(U\) a symmetric open neighbourhood of \(e\) in \(G\) and \(\psi: U^2 \to K^n\) a chart of \(G\) such that \(\psi(e) = 0\) . If \((x_1, \ldots, x_n)\) are the coordinates of \(x \in \psi(U)\) and \((y_1, \ldots, y_n)\) the coordinates of \(y \in \psi(U)\) , we denote by

\[
m _ {1} (x _ {1}, \dots , x _ {n}, y _ {1}, \dots , y _ {n}), \dots , m _ {n} (x _ {1}, \dots , x _ {n}, y _ {1}, \dots , y _ {n})
\]

the coordinates of \(\psi (\psi^{-1}(x)^{-1}\psi^{-1}(y))\) . Then, if we write, for \(1\leqslant k\leqslant n\)

\[
\begin{array}{r l} \varpi_ {k} (x _ {1}, \ldots , x _ {n}) = D _ {n + 1} m _ {k} (x _ {1}, \ldots , x _ {n}, x _ {1}, \ldots , x _ {n}) d x _ {1} + \dots \\ & \quad + \mathrm{D} _ {2 n} m _ {k} (x _ {1}, \ldots , x _ {n}, x _ {1}, \ldots , x _ {n}) d x _ {n}, \end{array}\tag{32}
\]

the differential forms \(\varpi_k\) on \(\psi(\mathbf{U})\) are derived through \(\psi\) from left invariant differential forms on \(\mathbf{G}\) and are such that \(\varpi_k(0,\ldots,0) = dx_k\) .

We apply Lemma 2 with E = K, taking \(\phi(g)\) to be the coordinate of \(\psi(g)\) of index k. We obtain a differential form \(\omega_{k}\) ; let \(\varpi_{k}\) be its transform under \(\psi\) . The value of \(\varpi_{k}\) at \((x_{1},\ldots,x_{n})\) is the differential at \((x_{1},\ldots,x_{n})\) of the function \(y\mapsto m_{k}(x_{1},\ldots,x_{n},y_{1},\ldots,y_{n})\) ; this value is therefore given by formula (32). It then suffices to use the conclusion to Lemma 2.

PROPOSITION 53. Let G be a Lie group, A a complete normable algebra and \(\phi\) a Lie group morphism of G into A*. For all \(g \in G\) , let \(\omega_g = \phi(g)^{-1}.d_g\phi\) . Then \(\omega\) is the left invariant differential form on G whose value at \(e\) is \(d_e\phi\) .

We apply Lemma 2 with \(\mathbf{E} = \mathbf{A}\) and \(\mathbf{U} = \mathbf{G}\) . The differential at \(g\) of the mapping \(h \mapsto \phi(g^{-1}h) = \phi(g)^{-1}\phi(h)\) is \(\phi(g)^{-1}.d_g\phi\) .

